\section{Zariski's Main Theorem}
\label{section-Zariski}

\noindent
In this section our aim is to prove the algebraic version of
Zariski's Main theorem. This theorem will be the basis of many
further developments in the theory of schemes and morphisms of
schemes later in the Stacks project.

\medskip\noindent
Let $R \to S$ be a ring map of finite type.
Our goal in this section is to show that the set
of points of $\Spec(S)$ where the map is quasi-finite
is {\it open} (Theorem \ref{theorem-main-theorem}).
In fact, it will turn out that there exists
a finite ring map $R \to S'$ such that in some sense the quasi-finite
locus of $S/R$ is open in $\Spec(S')$
(but we will not prove this in the algebra chapter since we do not
develop the language of schemes here -- for the case where $R \to S$
is quasi-finite see Lemma \ref{lemma-quasi-finite-open-integral-closure}).
These statements are somewhat tricky to prove and
we do it by a long list of lemmas concerning integral and
finite extensions of rings. This material may be found
in \cite{Henselian}, and \cite{Peskine}. We also found notes
by Thierry Coquand helpful.

\begin{lemma}
\label{lemma-make-integral-trivial}
Let $\varphi : R \to S$ be a ring map.
Suppose $t \in S$ satisfies the
relation $\varphi(a_0) + \varphi(a_1)t + \ldots + \varphi(a_n) t^n = 0$.
Then $\varphi(a_n)t$ is integral over $R$.
\end{lemma}

\begin{proof}
Namely, multiply the equation
$\varphi(a_0) + \varphi(a_1)t + \ldots + \varphi(a_n) t^n = 0$
with $\varphi(a_n)^{n-1}$ and write it as
$\varphi(a_0 a_n^{n-1}) +
\varphi(a_1 a_n^{n-2}) (\varphi(a_n)t) +
\ldots +
(\varphi(a_n) t)^n = 0$.
\end{proof}

\noindent
The following lemma is in some sense the key lemma in this
section.

\begin{lemma}
\label{lemma-make-integral-trick}
Let $R$ be a ring. Let $\varphi : R[x] \to S$ be
a ring map. Let $t \in S$.
Assume that (a) $t$ is integral over $R[x]$,
and (b) there exists a monic $p \in R[x]$ such that
$t \varphi(p) \in \Im(\varphi)$. Then there
exists a $q \in R[x]$ such that $t - \varphi(q)$
is integral over $R$.
\end{lemma}

\begin{proof}
Write $t \varphi(p) = \varphi(r)$ for some $r \in R[x]$.
Using euclidean division, write $r = qp + r'$ with
$q, r' \in R[x]$ and $\deg(r') < \deg(p)$. We may replace
$t$ by $t - \varphi(q)$ which is still integral over
$R[x]$, so that we obtain $t \varphi(p) = \varphi(r')$.
In the ring $S_t$ we may write this as
$\varphi(p) - (1/t) \varphi(r') = 0$.
This implies that $\varphi(x)$ gives an element of the
localization $S_t$ which is integral over
$\varphi(R)[1/t] \subset S_t$. On the other hand,
$t$ is integral over the subring $\varphi(R)[\varphi(x)] \subset S$.
Combined we conclude that $t$ is integral over
the subring $\varphi(R)[1/t] \subset S_t$, see Lemma
\ref{lemma-integral-transitive}. In other words
there exists an equation of the form
$$
t^d + \sum\nolimits_{i < d}
\left(\sum\nolimits_{j = 0, \ldots, n_i} \varphi(r_{i, j})/t^j\right) t^i = 0
$$
in $S_t$ with $r_{i, j} \in R$. This means that
$t^{d + N} +
\sum_{i < d} \sum_{j = 0, \ldots, n_i} \varphi(r_{i, j}) t^{i + N - j} = 0$
in $S$ for some $N$ large enough. In other words
$t$ is integral over $R$.
\end{proof}

\begin{lemma}
\label{lemma-combine-lemmas}
Let $R$ be a ring. Let $\varphi : R[x] \to S$ be
a ring map. Let $t \in S$. Assume $t$ is integral
over $R[x]$. Let $p \in R[x]$, $p = a_0 + a_1x + \ldots +
a_k x^k$ such that $t \varphi(p) \in \Im(\varphi)$.
Then there exists a $q \in R[x]$ and $n \geq 0$
such that $\varphi(a_k)^n t - \varphi(q)$ is integral
over $R$.
\end{lemma}

\begin{proof}
Let $R'$ and $S'$ be the localization of $R$ and $S$ at the element $a_k$.
Let $\varphi' : R'[x] \to S'$ be the localization of $\varphi$.
Let $t' \in S'$ be the image of $t$. Set $p' = p/a_k \in R'[x]$.
Then $t' \varphi'(p') \in \Im(\varphi')$ since $t \varphi(p) \in \Im(\varphi)$.
As $p'$ is monic, by
Lemma \ref{lemma-make-integral-trick} there exists a $q' \in R'[x]$
such that $t' - \varphi'(q')$ is integral over $R'$.
We may choose an $n \geq 0$ and an element $q \in R[x]$
such that $a_k^n q'$ is the image of $q$.
Then $\varphi(a_k)^n t - \varphi(q)$ is an element of $S$ whose
image in $S'$ is integral over $R'$.
By Lemma \ref{lemma-integral-closure-localize}
there exists an $m \geq 0$ such that
$\varphi(a_k)^m(\varphi(a_k)^n t - \varphi(q))$ is integral over $R$.
Thus $\varphi(a_k)^{m + n}t - \varphi(a_k^m q)$
is integral over $R$ as desired.
\end{proof}

\begin{situation}
\label{situation-one-transcendental-element}
Let $R$ be a ring.
Let $\varphi : R[x] \to S$ be finite.
Let
$$
J = \{ g \in S \mid gS \subset \Im(\varphi)\}
$$
be the ``conductor ideal'' of $\varphi$.
Assume $\varphi(R) \subset S$ integrally closed in $S$.
\end{situation}

\begin{lemma}
\label{lemma-leading-coefficient-in-J}
In Situation \ref{situation-one-transcendental-element}.
Suppose $u \in S$, $a_0, \ldots, a_k \in R$,
$u \varphi(a_0 + a_1x + \ldots + a_k x^k) \in J$.
Then there exists an $m \geq 0$ such that
$u \varphi(a_k)^m \in J$.
\end{lemma}

\begin{proof}
Assume that $S$ is generated by $t_1, \ldots, t_n$
as an $R[x]$-module. In this case
$J = \{ g \in S \mid gt_i \in \Im(\varphi)\text{ for all }i\}$.
Note that each element $u t_i$ is integral over
$R[x]$, see Lemma \ref{lemma-finite-is-integral}.
We have $\varphi(a_0 + a_1x + \ldots + a_k x^k) u t_i \in
\Im(\varphi)$. By Lemma \ref{lemma-combine-lemmas}, for
each $i$ there exists an integer $n_i$ and an element
$q_i \in R[x]$ such that $\varphi(a_k^{n_i}) u t_i - \varphi(q_i)$
is integral over $R$. By assumption this element is in $\varphi(R)$
and hence $\varphi(a_k^{n_i}) u t_i \in \Im(\varphi)$.
It follows that $m = \max\{n_1, \ldots, n_n\}$ works.
\end{proof}

\begin{lemma}
\label{lemma-all-coefficients-in-J}
In Situation \ref{situation-one-transcendental-element}.
Suppose $u \in S$, $a_0, \ldots, a_k \in R$,
$u \varphi(a_0 + a_1x + \ldots + a_k x^k) \in \sqrt{J}$.
Then $u \varphi(a_i) \in \sqrt{J}$ for all $i$.
\end{lemma}

\begin{proof}
Under the assumptions of the lemma we have
$u^n \varphi(a_0 + a_1x + \ldots + a_k x^k)^n \in J$ for
some $n \geq 1$. By Lemma \ref{lemma-leading-coefficient-in-J}
we deduce $u^n \varphi(a_k^{nm}) \in J$ for some $m \geq 1$.
Thus $u \varphi(a_k) \in \sqrt{J}$, and so
$u \varphi(a_0 + a_1x + \ldots + a_k x^k) - u \varphi(a_k x^k) =
u \varphi(a_0 + a_1x + \ldots + a_{k-1} x^{k-1}) \in \sqrt{J}$.
We win by induction on $k$.
\end{proof}

\noindent
This lemma suggests the following definition.

\begin{definition}
\label{definition-strongly-transcendental}
Given an inclusion of rings $R \subset S$ and
an element $x \in S$ we say that $x$ is
{\it strongly transcendental over $R$} if
whenever $u(a_0 + a_1 x + \ldots + a_k x^k) = 0$
with $u \in S$ and $a_i \in R$, then
we have $ua_i = 0$ for all $i$.
\end{definition}

\noindent
Note that if $S$ is a domain then this is the same as
saying that $x$ as an element of the fraction field of
$S$ is transcendental over the fraction field of $R$.

\begin{lemma}
\label{lemma-reduced-strongly-transcendental-minimal-prime}
Suppose $R \subset S$ is an inclusion of reduced rings
and suppose that $x \in S$ is strongly transcendental over $R$.
Let $\mathfrak q \subset S$ be a minimal prime
and let $\mathfrak p = R \cap \mathfrak q$.
Then the image of $x$ in $S/\mathfrak q$ is strongly
transcendental over the subring $R/\mathfrak p$.
\end{lemma}

\begin{proof}
Suppose $u(a_0 + a_1x + \ldots + a_k x^k) \in \mathfrak q$.
By Lemma \ref{lemma-minimal-prime-reduced-ring}
the local ring $S_{\mathfrak q}$ is a field,
and hence $u(a_0 + a_1x + \ldots + a_k x^k) $ is zero
in $S_{\mathfrak q}$. Thus $uu'(a_0 + a_1x + \ldots + a_k x^k) = 0$
for some $u' \in S$, $u' \not\in \mathfrak q$.
Since $x$ is strongly transcendental over $R$ we get
$uu'a_i = 0$ for all $i$. This in turn implies
that $ua_i \in \mathfrak q$.
\end{proof}

\begin{lemma}
\label{lemma-domains-transcendental-not-quasi-finite}
Suppose $R\subset S$ is an inclusion of domains and
let $x \in S$. Assume $x$ is (strongly) transcendental over $R$
and that $S$ is finite over $R[x]$. Then $R \to S$ is not
quasi-finite at any prime of $S$.
\end{lemma}

\begin{proof}
As a first case, assume that $R$ is normal, see
Definition \ref{definition-ring-normal}.
By Lemma \ref{lemma-polynomial-ring-normal}
we see that $R[x]$ is normal.
Take a prime $\mathfrak q \subset S$,
and set $\mathfrak p = R \cap \mathfrak q$.
Assume that the extension $\kappa(\mathfrak p)
\subset \kappa(\mathfrak q)$ is finite.
This would be the case if $R \to S$ is
quasi-finite at $\mathfrak q$.
Let $\mathfrak r = R[x] \cap \mathfrak q$.
Then since $\kappa(\mathfrak p)
\subset \kappa(\mathfrak r) \subset \kappa(\mathfrak q)$
we see that the extension $\kappa(\mathfrak p)
\subset \kappa(\mathfrak r)$ is finite too.
Thus the inclusion $\mathfrak r \supset \mathfrak p R[x]$
is strict. By going down for $R[x] \subset S$,
see Proposition \ref{proposition-going-down-normal-integral},
we find a prime $\mathfrak q' \subset \mathfrak q$,
lying over the prime $\mathfrak pR[x]$. Hence
the fibre $\Spec(S \otimes_R \kappa(\mathfrak p))$
contains a point not equal to $\mathfrak q$,
namely $\mathfrak q'$, whose closure contains $\mathfrak q$ and hence
$\mathfrak q$ is not isolated in its fibre.

\medskip\noindent
If $R$ is not normal, let $R \subset R' \subset K$ be
the integral closure $R'$ of $R$ in its field of fractions
$K$. Let $S \subset S' \subset L$ be the subring $S'$ of
the field of fractions $L$ of $S$ generated by $R'$ and
$S$. Note that by construction the map $S \otimes_R R'
\to S'$ is surjective. This implies that $R'[x] \subset S'$
is finite. Also, the map $S \subset S'$
induces a surjection on $\Spec$, see
Lemma \ref{lemma-integral-overring-surjective}.
We conclude by Lemma \ref{lemma-four-rings} and the normal case
we just discussed.
\end{proof}

\begin{lemma}
\label{lemma-reduced-strongly-transcendental-not-quasi-finite}
Suppose $R \subset S$ is an inclusion of reduced rings.
Assume $x \in S$ be strongly transcendental over $R$,
and $S$ finite over $R[x]$. Then $R \to S$ is not
quasi-finite at any prime of $S$.
\end{lemma}

\begin{proof}
Let $\mathfrak q \subset S$ be any prime.
Choose a minimal prime $\mathfrak q' \subset \mathfrak q$.
According to Lemmas
\ref{lemma-reduced-strongly-transcendental-minimal-prime} and
\ref{lemma-domains-transcendental-not-quasi-finite}
the extension $R/(R \cap \mathfrak q') \subset
S/\mathfrak q'$ is not quasi-finite at the prime corresponding
to $\mathfrak q$. By Lemma \ref{lemma-four-rings}
the extension $R \to S$ is not quasi-finite
at $\mathfrak q$.
\end{proof}

\begin{lemma}
\label{lemma-quasi-finite-monogenic}
Let $R$ be a ring. Let $S = R[x]/I$.
Let $\mathfrak q \subset S$ be a prime.
Assume $R \to S$ is quasi-finite at $\mathfrak q$.
Let $S' \subset S$ be the integral closure of $R$ in $S$.
Then there exists an element
$g \in S'$, $g \not\in \mathfrak q$ such that
$S'_g \cong S_g$.
\end{lemma}

\begin{proof}
Let $\mathfrak p$ be the image of $\mathfrak q$ in $\Spec(R)$.
There exists an $f \in I$, $f = a_nx^n + \ldots + a_0$ such that
$a_i \not \in \mathfrak p$ for some $i$. Namely, otherwise the fibre ring
$S \otimes_R \kappa(\mathfrak p)$ would be $\kappa(\mathfrak p)[x]$
and the map would not be quasi-finite at any prime lying
over $\mathfrak p$. We conclude there exists a relation
$b_m x^m + \ldots + b_0 = 0$ with $b_j \in S'$, $j = 0, \ldots, m$
and $b_j \not \in \mathfrak q \cap S'$ for some $j$.
We prove the lemma by induction on $m$. The base case is
$m = 0$ is vacuous (because the statements $b_0 = 0$ and
$b_0 \not \in \mathfrak q$ are contradictory).

\medskip\noindent
The case $b_m \not \in \mathfrak q$. In this case $x$ is integral
over $S'_{b_m}$, in fact $b_mx \in S'$ by
Lemma \ref{lemma-make-integral-trivial}.
Hence the injective map $S'_{b_m} \to S_{b_m}$ is also surjective, i.e.,
an isomorphism as desired.

\medskip\noindent
The case $b_m \in \mathfrak q$. In this case we have $b_mx \in S'$ by
Lemma \ref{lemma-make-integral-trivial}.
Set $b'_{m - 1} = b_mx + b_{m - 1}$. Then
$$
b'_{m - 1}x^{m - 1} + b_{m - 2}x^{m - 2} + \ldots + b_0 = 0
$$
Since $b'_{m - 1}$ is congruent to $b_{m - 1}$ modulo $S' \cap \mathfrak q$
we see that it is still the case that one of
$b'_{m - 1}, b_{m - 2}, \ldots, b_0$ is not in $S' \cap \mathfrak q$.
Thus we win by induction on $m$.
\end{proof}

\begin{theorem}[Zariski's Main Theorem]
\label{theorem-main-theorem}
Let $R$ be a ring. Let $R \to S$ be a finite type $R$-algebra.
Let $S' \subset S$ be the integral closure of $R$ in $S$.
Let $\mathfrak q \subset S$ be a prime of $S$.
If $R \to S$ is quasi-finite at $\mathfrak q$ then
there exists a $g \in S'$, $g \not \in \mathfrak q$
such that $S'_g \cong S_g$.
\end{theorem}

\begin{proof}
There exist finitely many elements
$x_1, \ldots, x_n \in S$ such that $S$ is finite
over the $R$-sub algebra generated by $x_1, \ldots, x_n$. (For
example generators of $S$ over $R$.) We prove the proposition
by induction on the minimal such number $n$.

\medskip\noindent
The case $n = 0$ is trivial, because in this case $S' = S$,
see Lemma \ref{lemma-finite-is-integral}.

\medskip\noindent
The case $n = 1$. We may replace $R$ by its integral closure in $S$
(Lemma \ref{lemma-quasi-finite-permanence} guarantees that $R \to S$
is still quasi-finite at $\mathfrak q$). Thus we may assume
$R \subset S$ is integrally closed in $S$, in other words $R = S'$.
Consider the map $\varphi : R[x] \to S$, $x \mapsto x_1$.
(We will see that $\varphi$ is not injective below.)
By assumption $\varphi$ is finite. Hence we are in Situation
\ref{situation-one-transcendental-element}.
Let $J \subset S$ be the ``conductor ideal'' defined
in Situation \ref{situation-one-transcendental-element}.
Consider the diagram
$$
\xymatrix{
R[x] \ar[r] & S \ar[r] & S/\sqrt{J} & R/(R \cap \sqrt{J})[x] \ar[l]
\\
& R \ar[lu] \ar[r] \ar[u] & R/(R \cap \sqrt{J}) \ar[u] \ar[ru] &
}
$$
According to Lemma \ref{lemma-all-coefficients-in-J}
the image of $x$ in the quotient $S/\sqrt{J}$
is strongly transcendental over $R/ (R \cap \sqrt{J})$.
Hence by Lemma \ref{lemma-reduced-strongly-transcendental-not-quasi-finite}
the ring map $R/ (R \cap \sqrt{J}) \to S/\sqrt{J}$
is not quasi-finite at any prime of $S/\sqrt{J}$.
By Lemma \ref{lemma-four-rings} we deduce that $\mathfrak q$
does not lie in $V(J) \subset \Spec(S)$.
Thus there exists an element $s \in J$,
$s \not\in \mathfrak q$. By definition of $J$ we may write
$s = \varphi(f)$ for some polynomial $f \in R[x]$.
Let $I = \Ker(\varphi : R[x] \to S)$. Since $\varphi(f) \in J$
we get $(R[x]/I)_f \cong S_{\varphi(f)}$. Also $s \not \in \mathfrak q$
means that $f \not \in \varphi^{-1}(\mathfrak q)$. Thus
$\varphi^{-1}(\mathfrak q)/I$ is a prime of $R[x]/I$
at which $R \to R[x]/I$ is quasi-finite, see
Lemma \ref{lemma-quasi-finite-local}. Note that $R$ is integrally closed
in $R[x]/I$ since $R$ is integrally closed in $S$. By
Lemma \ref{lemma-quasi-finite-monogenic}
there exists an element $h \in R$, $h \not \in R \cap \mathfrak q$
such that $R_h \cong (R[x]/I)_h$. Thus
$(R[x]/I)_{fh} = S_{\varphi(fh)}$ is isomorphic to a principal
localization $R_{h'}$ of $R$ for some
$h' \in R$, $h' \not \in \mathfrak q$.

\medskip\noindent
The case $n > 1$. Consider the subring $R' \subset S$
which is the integral closure of $R[x_1, \ldots, x_{n-1}]$
in $S$. By Lemma \ref{lemma-quasi-finite-permanence} the extension
$S/R'$ is quasi-finite at $\mathfrak q$. Also, note
that $S$ is finite over $R'[x_n]$.
By the case $n = 1$ above, there exists a $g' \in R'$,
$g' \not \in \mathfrak q$ such that
$(R')_{g'} \cong S_{g'}$. At this point we cannot
apply induction to $R \to R'$ since $R'$ may not be finite type over $R$.
Since $S$ is finitely generated over $R$ we deduce in particular
that $(R')_{g'}$ is finitely generated over $R$. Say
the elements $g'$, and $y_1/(g')^{n_1}, \ldots, y_N/(g')^{n_N}$
with $y_i \in R'$
generate $(R')_{g'}$ over $R$. Let $R''$ be the $R$-sub algebra
of $R'$ generated by $x_1, \ldots, x_{n-1}, y_1, \ldots, y_N, g'$.
This has the property $(R'')_{g'} \cong S_{g'}$. Surjectivity
because of how we chose $y_i$, injectivity because
$R'' \subset R'$, and localization is exact. Note that
$R''$ is finite over $R[x_1, \ldots, x_{n-1}]$ because
of our choice of $R'$, see Lemma \ref{lemma-characterize-integral}.
Let $\mathfrak q'' = R'' \cap \mathfrak q$. Since
$(R'')_{\mathfrak q''} = S_{\mathfrak q}$ we see that
$R \to R''$ is quasi-finite at $\mathfrak q''$, see
Lemma \ref{lemma-isolated-point-fibre}.
We apply our induction hypothesis to $R \to R''$, $\mathfrak q''$
and $x_1, \ldots, x_{n-1} \in R''$ and we find a subring
$R''' \subset R''$ which is integral over $R$ and an
element $g'' \in R'''$, $g'' \not \in \mathfrak q''$
such that $(R''')_{g''} \cong (R'')_{g''}$. Write the image of $g'$ in
$(R'')_{g''}$ as $g'''/(g'')^n$ for some $g''' \in  R'''$.
Set $g = g''g''' \in R'''$. Then it is clear that $g \not\in
\mathfrak q$ and $(R''')_g \cong S_g$. Since by construction
we have $R''' \subset S'$ we also have $S'_g \cong S_g$ as desired.
\end{proof}

\begin{lemma}
\label{lemma-quasi-finite-open}
Let $R \to S$ be a finite type ring map.
The set of points $\mathfrak q$ of $\Spec(S)$ at which
$S/R$ is quasi-finite is open in $\Spec(S)$.
\end{lemma}

\begin{proof}
Let $\mathfrak q \subset S$ be a point at which the ring map
is quasi-finite. By Theorem \ref{theorem-main-theorem}
there exists an integral ring map $R \to S'$, $S' \subset S$
and an element $g \in S'$, $g\not \in \mathfrak q$ such that
$S'_g \cong S_g$. Since $S$ and hence $S_g$ are of finite type
over $R$ we may find finitely many elements
$y_1, \ldots, y_N$ of $S'$ such that $S''_g \cong S_g$
where $S'' \subset S'$ is the sub $R$-algebra generated
by $g, y_1, \ldots, y_N$. Since $S''$ is finite over $R$
(see Lemma \ref{lemma-characterize-integral}) we see that
$S''$ is quasi-finite over $R$ (see Lemma \ref{lemma-quasi-finite}).
It is easy to see that this implies that $S''_g$ is quasi-finite over $R$,
for example because the property of being quasi-finite at a prime depends
only on the local ring at the prime. Thus we see that $S_g$ is quasi-finite
over $R$. By the same token this implies that $R \to S$ is quasi-finite
at every prime of $S$ which lies in $D(g)$.
\end{proof}

\begin{lemma}
\label{lemma-quasi-finite-open-integral-closure}
Let $R \to S$ be a finite type ring map.
Suppose that $S$ is quasi-finite over $R$.
Let $S' \subset S$ be the integral closure of $R$ in $S$. Then
\begin{enumerate}
\item $\Spec(S) \to \Spec(S')$ is a homeomorphism
onto an open subset,
\item if $g \in S'$ and $D(g)$ is contained in the image
of the map, then $S'_g \cong S_g$, and
\item there exists a finite $R$-algebra $S'' \subset S'$
such that (1) and (2) hold for the ring map
$S'' \to S$.
\end{enumerate}
\end{lemma}

\begin{proof}
Because $S/R$ is quasi-finite we may apply
Theorem \ref{theorem-main-theorem} to
each point $\mathfrak q$ of $\Spec(S)$.
Since $\Spec(S)$ is quasi-compact, see
Lemma \ref{lemma-quasi-compact}, we may choose
a finite number of $g_i \in S'$, $i = 1, \ldots, n$
such that $S'_{g_i} = S_{g_i}$, and such that
$g_1, \ldots, g_n$ generate the unit ideal in $S$
(in other words the standard opens of $\Spec(S)$ associated
to $g_1, \ldots, g_n$ cover all of $\Spec(S)$).

\medskip\noindent
Suppose that $D(g) \subset \Spec(S')$
is contained in the image. Then $D(g) \subset \bigcup D(g_i)$.
In other words, $g_1, \ldots, g_n$ generate the unit ideal of
$S'_g$. Note that $S'_{gg_i} \cong S_{gg_i}$ by our choice
of $g_i$. Hence $S'_g \cong S_g$ by Lemma \ref{lemma-cover}.

\medskip\noindent
We construct a finite algebra $S'' \subset S'$ as
in (3). To do this note that each $S'_{g_i} \cong S_{g_i}$
is a finite type $R$-algebra. For each $i$ pick
some elements $y_{ij} \in S'$ such that each
$S'_{g_i}$ is generated as $R$-algebra by $1/g_i$
and the elements $y_{ij}$. Then set $S''$
equal to the sub $R$-algebra of $S'$ generated by all $g_i$
and all the $y_{ij}$. Details omitted.
\end{proof}




